\begin{thebibliography}{99}
\bibitem{dysthe1979}
K. B. Dysthe, \emph{Note on a modification to the nonlinear
Schr\"odinger equation for application to deep water waves},
Proc. R. Soc. Lond. A \textbf{369} (1979), 105--114.
Supplied paper, especially printed pp.~109--113.
\bibitem{schober}
C. M. Schober and A. Islas,
\emph{Higher-Order Extensions of Weakly Viscous Dysthe Theory and a
Phase-Lag Model for Nonlinear Mean-Flow Damping},
\href{https://arxiv.org/abs/2608.04473v1}{arXiv:2608.04473v1} (2026).
Supplied preprint, especially \S\S4--7.
\bibitem{fedele}
F. Fedele and D. Dutykh,
\emph{Hamiltonian description and traveling waves of the spatial Dysthe
equations}, \href{https://arxiv.org/abs/1110.3605v3}{arXiv:1110.3605v3}
(2012). Equation (1.3), \S\S3--5.
\bibitem{lomei}
E. Lo and C. C. Mei,
\emph{A numerical study of water-wave modulation based on a higher-order
nonlinear Schr\"odinger equation}, J. Fluid Mech. \textbf{150} (1985),
395--416.
\href{https://doi.org/10.1017/S0022112085000180}{doi:10.1017/S0022112085000180}.
\bibitem{recurrence}
A. Armaroli, M. Brunetti and J. Kasparian,
\emph{Recurrence in the high-order nonlinear Schr\"odinger equation:
a low dimensional analysis}, Phys. Rev. E \textbf{96} (2017), 012222.
\href{https://arxiv.org/abs/1703.09482v2}{arXiv:1703.09482v2},
\S\S II--III.
\bibitem{broadband}
K. Trulsen and K. B. Dysthe,
\emph{A modified nonlinear Schr\"odinger equation for broader bandwidth
gravity waves on deep water}, Wave Motion \textbf{24} (1996), 281--289.
\href{https://doi.org/10.1016/S0165-2125(96)00020-0}{doi:10.1016/S0165-2125(96)00020-0}.
\bibitem{viscous}
J. D. Carter and A. Govan,
\emph{Frequency downshift in a viscous fluid},
\href{https://arxiv.org/abs/1601.03932}{arXiv:1601.03932}.
Journal version (2016):
\href{https://doi.org/10.1016/j.euromechflu.2016.06.002}{doi:10.1016/j.euromechflu.2016.06.002}.
\bibitem{experiments}
J. D. Carter, D. Henderson and I. Butterfield,
\emph{A Comparison of Frequency Downshift Models of Wave Trains on Deep Water},
\href{https://arxiv.org/abs/1809.09536v2}{arXiv:1809.09536v2} (2018),
\S\S4--6.
\bibitem{transverse}
K. Trulsen, C. T. Stansberg and K. B. Dysthe,
\emph{Experimental evidence of transverse modulation and frequency
downshift of uniform waves in a long tank},
J. Fluid Mech. \textbf{1024} (2025), A14.
\href{https://doi.org/10.1017/jfm.2025.10822}{doi:10.1017/jfm.2025.10822},
\S\S5--6.
\bibitem{fno}
Z. Li et al.,
\emph{Fourier Neural Operator for Parametric Partial Differential Equations},
ICLR (2021).
\href{https://arxiv.org/abs/2010.08895}{arXiv:2010.08895}, \S\S3--5.
\bibitem{pino}
Z. Li et al.,
\emph{Physics-Informed Neural Operator for Learning Partial Differential
Equations},
\href{https://arxiv.org/abs/2111.03794v4}{arXiv:2111.03794v4},
\S3 and Appendix E.
\bibitem{deeponet}
L. Lu, P. Jin and G. E. Karniadakis,
\emph{DeepONet: Learning nonlinear operators for identifying differential
equations based on the universal approximation theorem of operators},
\href{https://arxiv.org/abs/1910.03193v3}{arXiv:1910.03193v3} (2020),
\S2. This citation identifies the inspected preprint version.
\bibitem{pideeponet}
S. Wang, H. Wang and P. Perdikaris,
\emph{Learning the solution operator of parametric partial differential
equations with physics-informed DeepONets},
Science Advances \textbf{7} (2021), eabi8605.
\href{https://pmc.ncbi.nlm.nih.gov/articles/PMC8480920/}{Open primary article}.
\bibitem{pushforward}
J. Brandstetter, D. Worrall and M. Welling,
\emph{Message Passing Neural PDE Solvers}, ICLR (2022).
\href{https://arxiv.org/abs/2202.03376v2}{arXiv:2202.03376v2}, \S3.1.
\bibitem{stability}
M. McCabe, P. Harrington, S. Subramanian and J. Brown,
\emph{Towards Stability of Autoregressive Neural Operators},
Transactions on Machine Learning Research (2023).
\href{https://arxiv.org/abs/2306.10619v2}{arXiv:2306.10619v2}, \S\S3--4.
\bibitem{causal}
S. Wang, S. Sankaran and P. Perdikaris,
\emph{Respecting causality is all you need for training physics-informed
neural networks},
\href{https://arxiv.org/abs/2203.07404}{arXiv:2203.07404}, \S\S3--4.
Journal version, Comput. Methods Appl. Mech. Eng. (2024):
\href{https://doi.org/10.1016/j.cma.2024.116813}{doi:10.1016/j.cma.2024.116813}.
\bibitem{failures}
A. Krishnapriyan et al.,
\emph{Characterizing possible failure modes in physics-informed neural
networks}, NeurIPS (2021).
\href{https://arxiv.org/abs/2109.01050}{arXiv:2109.01050}.
\bibitem{mps}
R. Mosincat, D. Pilod and J.-C. Saut,
\emph{Global well-posedness and scattering for the Dysthe equation
in $L^2(\RR^2)$},
\href{https://arxiv.org/abs/2007.01613}{arXiv:2007.01613}.
Theorems 1.1--1.2 concern a different temporal, whole-space formulation.
\bibitem{mishra}
S. Mishra and R. Molinaro,
\emph{Estimates on the generalization error of Physics Informed Neural
Networks (PINNs) for approximating PDEs},
\href{https://arxiv.org/abs/2006.16144}{arXiv:2006.16144}.
Used for the need for stability and sampling hypotheses, not a Dysthe theorem.
\bibitem{ehlers2024}
S. Ehlers, N. A. Wagner, A. Scherzl, M. Klein, N. Hoffmann and M. Stender,
\emph{Data Assimilation and Parameter Identification for Water Waves
Using the Nonlinear Schr\"odinger Equation and Physics-Informed Neural
Networks}, Fluids \textbf{9} (2024), 231.
\href{https://doi.org/10.3390/fluids9100231}{doi:10.3390/fluids9100231},
\S\S2.3--4.
\bibitem{ehlers2025}
S. Ehlers, M. Stender and N. Hoffmann,
\emph{Bridging ocean wave physics and deep learning: Physics-informed
neural operators for nonlinear wavefield reconstruction in real-time},
\href{https://arxiv.org/abs/2508.03315v1}{arXiv:2508.03315v1}
(2025, preprint), \S II.3.
\bibitem{eeltink}
D. Eeltink et al., \emph{Nonlinear wave evolution with data-driven breaking},
Nature Communications \textbf{13} (2022), 2343.
\href{https://doi.org/10.1038/s41467-022-30025-z}{doi:10.1038/s41467-022-30025-z},
Methods and Discussion.
\end{thebibliography}
